\label{section:vector-register-model}
The scalable-vector extension provides 32 vector registers, \texttt{V0}
through \texttt{V31}, and 16 predicate registers, \texttt{P0} through
\texttt{P15}.  A vector register contains VLEN bits.  VLEN is an
implementation-selected power of two from 128 through 2048 bits and remains
stable until reset.  A predicate register contains one bit for each byte of a
vector register; its packed image therefore occupies VLEN/64 bytes.  Reset
clears every vector and predicate bit.

The \texttt{VECTOR} discovery predicate and the \texttt{VECTOR\_PARAMETERS}
CPUID leaf report the extension and VLEN.  Integer and raw-bit vector forms
require \texttt{VECTOR}.  A form whose selected element type is \texttt{H},
\texttt{S}, or \texttt{D} additionally requires the base floating-point
extension.  The requirement follows the decoded element type, including a
shared integer/floating-point opcode; disabling FP does not disable the
integer forms of that opcode.  Vector floating-point operations use the
ordinary \texttt{FSTATUS} environment and accrue active-lane causes into
\texttt{FFLAGS}.

\subsection{Element Types and Assembly Spelling}

A size-bearing vector instruction has one suffix.  The integer element
suffixes \texttt{B}, \texttt{W}, \texttt{L}, and \texttt{Q} select 1-, 2-,
4-, and 8-byte elements.  The floating-point suffixes \texttt{H}, \texttt{S},
and \texttt{D} select IEEE 754 binary16, binary32, and binary64 elements.
Binary16 is part of the vector FP facility; its default quiet NaN is
\texttt{0x7e00}.

An opcode family that supports both integer and floating-point arithmetic
uses the three-bit \texttt{tzz} selector.  Values zero through three select
\texttt{B/W/L/Q}, value four is invalid, and values five through seven select
\texttt{H/S/D}.  A family whose numeric domain is fixed by its opcode uses a
two-bit \texttt{zz} selector: integer families assign all four values to
\texttt{B/W/L/Q}, while FP families reserve zero and assign one through three
to \texttt{H/S/D}.  A reserved selector does not name another data format.

The suffix is the source width for a width-changing instruction; the
destination width and signedness are part of the mnemonic.  For example,
\texttt{VEXTZQ.B} zero-extends a byte into a quadword.  An instruction whose
suffix controls only lane width accepts \texttt{H/S/D} as assembly aliases for
\texttt{W/L/Q}; canonical disassembly uses \texttt{B/W/L/Q}.  Predicate
operations, \texttt{RDCNT}, \texttt{PLOOP}, \texttt{VMOV}, \texttt{VMOVZ},
and the permute, slide, slice, and zip families use this width-only rule.

\subsection{Predicates and Destination Images}

For an element size of $E$ bytes, logical lane $i$ is governed by predicate
bit $iE$.  All other predicate bits are nonsignificant for that operation.
Unless an instruction explicitly produces a complete predicate image or uses
zeroing movement, a false governing bit annuls that lane: it performs no
memory access, produces no exception, and preserves the old destination lane.

An instruction that writes a predicate result writes a complete packed
predicate image.  It sets or clears each significant result bit and clears
all nonsignificant bits.  Integer comparisons support \texttt{EQ},
\texttt{NE}, \texttt{ULT}, \texttt{UGE}, \texttt{ULE}, \texttt{UGT},
\texttt{LT}, \texttt{GE}, \texttt{LE}, and \texttt{GT}.  FP comparisons
support \texttt{EQ}, \texttt{NE}, \texttt{LT}, \texttt{LE}, \texttt{GE},
\texttt{GT}, \texttt{VS} for unordered, and \texttt{VC} for ordered.
\texttt{VTESTZ} and \texttt{VTESTNZ} are distinct integer test operations.
Vector integer operations do not update scalar \texttt{FLAGS}.

\texttt{PLOOP} consumes a remaining-count register and an offset register.  If
the remaining count is zero, it takes its encoded branch target and does not
change the predicate or either register.  Otherwise, an offset greater than
or equal to the selected lane count raises \texttt{VECTOR\_RANGE\_ERROR} with
error code zero and leaves the instruction uncommitted.  A valid nonzero
iteration activates the next
$\min(remaining, lane\_count-offset)$ lanes, subtracts that number from the
remaining count, clears the offset register, and falls through.  Software
advances the offset between chunks when it needs a nonzero subsequent chunk
origin.

\subsection{Lane Mapping and Scalar Bridges}

Vector bytes use increasing-address little-endian order.  A same-width
operation treats a vector as VLEN divided by the selected element width lanes.
For widening, narrowing, FP format conversion, and integer/FP conversion, the
operation container is the larger of the source and destination widths.  Each
active container reads the source value from its low source-width part and
writes the result into its low destination-width part.  Other bits of an
active destination container and every inactive container remain unchanged.

\texttt{VDUP} broadcasts a scalar register, FP register, or immediate into
every lane.  \texttt{VEXTRACT} and \texttt{VINSERT} use a scalar lane index.  An
index at or above the lane count raises \texttt{VECTOR\_RANGE\_ERROR} with
error code one and commits no state.  \texttt{RDVL} returns VLEN in bytes;
\texttt{RDCNT} returns the selected lane count.  \texttt{ADDVL} and
\texttt{ADDPL} scale their signed immediate by VLEN bytes and packed predicate
bytes, respectively.

\subsection{Vector Memory Operations}

A vector memory operand uses the ordinary compact, \texttt{EXT1}, and
\texttt{EXT2} descriptor formats in vector context.  The descriptor computes
one scalar anchor address.  The vector-context compact table supplies
\texttt{VSTRIDE} selectors where the scalar table has stack-relative and
immediate forms.  A contiguous access uses

\[
  address[i] = anchor + iE,
\]

and \texttt{VSTRIDE} uses

\[
  address[i] = base + i\,signed64(stride) + displacement.
\]

All arithmetic is modulo $2^{64}$.  Existing effective-address auto-update is
applied once to the scalar anchor in the ordinary architectural order; it is
not repeated per lane.

Inactive lanes issue no memory transaction.  \texttt{VMOV} merges a memory
source into its destination, while \texttt{VMOVZ} writes zero to inactive
destination lanes.  A memory-destination arithmetic form is a lane-wise
read-modify-write, not an atomic operation.  A vector instruction nevertheless
has one architectural commit point: address evaluation, all active reads,
computation, and all active writes must succeed before any register,
predicate, memory, auto-update, \texttt{FFLAGS}, or PC effect becomes visible.
If more than one active lane faults, the implementation may select any one of
those faults, but the reported fault address must identify the selected lane.

\subsection{Floating-Point Exceptions, Conversions, and Reductions}

Each active FP lane applies the corresponding scalar rounding, NaN,
signed-zero, default-NaN, denormal, and exception rules.  Inactive lanes
generate no cause.  The instruction unions all active-lane causes before
testing the enables in the original \texttt{FSTATUS}.  If any generated cause
is enabled, one \texttt{FLOATING\_POINT\_FAULT} occurs and neither a destination
effect nor newly accrued \texttt{FFLAGS} becomes visible.  Otherwise the union
is accrued once and all destination effects commit together.  FP-to-integer
conversion rounds toward zero and follows the scalar saturation and invalid
result rules.

Integer reductions return a zero-extended result in an \texttt{Rn}.  For an
empty predicate, ADD, OR, and XOR return zero; AND and unsigned MIN return the
all-ones value of the selected width; signed MIN returns the signed maximum;
signed MAX returns the signed minimum; and unsigned MAX returns zero.  FP
reductions return an \texttt{Fn}.  Empty FP ADD returns positive zero, empty FP
MIN returns positive infinity, and empty FP MAX returns negative infinity.
FP ADD visits active lanes in increasing logical-lane order and may not be
reassociated.

\subsection{Execution and Saved State}

The vector SAVE/RESTORE component contains the complete scalable vector and
packed predicate images.  A committed \texttt{Vn} or \texttt{Pn} write marks
that component modified.  Architectural event entry does not otherwise change
vector state.

An implementation may internally coalesce eligible scalar \texttt{REP} or
\texttt{REPcc} iterations into vector-sized work.  This permission creates no
vector instruction, predicate, or additional repeat state.  The observable
counter, condition, \texttt{FLAGS}, memory, auto-update, fault, event, and
restart behavior must equal the scalar iteration sequence, and only the
scalar prefix preceding the first fault or terminating \texttt{REPcc}
observation may commit.
